\section{Introduction}
Verdant-Minds explores whether structured memory and dynamic state regulation can yield measurable emergent concept organization. The project operates across a frozen v1 substrate and an active v2 research stack. Our central empirical focus is a temporal pattern in the emergent subgraph: emergent concepts added later tend to connect backward in time to earlier emergent concepts.

\subsection{Contributions}
This paper makes five contributions:
\begin{enumerate}[leftmargin=*]
  \item \textbf{Hybrid architecture}: a graph-memory substrate, ECWF wave dynamics, a governance layer (``Three Kings''), and a nine-stage processing pipeline.
  \item \textbf{Instrumentation}: coherence invariants (HCI, triangle validity, and alpha-critical) plus thermodynamic phase tracking via $T_g$.
  \item \textbf{Emergence mechanism}: an operational emergence detector combining entropy windowing, magnitude thresholding, and novelty scoring.
  \item \textbf{Empirical finding}: emergent$\rightarrow$emergent edges exhibit strong backward temporal bias, indicating preferential attachment from newer emergent concepts to older emergent concepts.
  \item \textbf{Replication}: 20 seed-controlled runs with 20/20 runs at $z>2$ under shuffle null.
\end{enumerate}
